% !TEX root = ../main.tex

\chapter{数据集构建与模型训练方法}

\section{基础模型与任务形式定义}

本文以 Qwen2.5-7B-Instruct 作为基础模型，在此基础上做监督微调，并预留后续强化学习训练流程。选择指令模型的主要原因是，它已经能够较好地理解中文任务说明，也能按照给定格式回答。对于本文任务来说，模型需要同时阅读 K 线描述、新闻摘要和研报摘要，再输出方向判断或趋势分析，因此从指令模型开始会比从普通语言模型开始更方便。

股票预测没有被处理为精确涨跌幅回归，而是被定义为方向分类。单一预测周期只包含“涨”和“跌”两类输出；多周期预测则要求模型同时给出未来 1 个、3 个和 7 个交易日后的方向结果。这样会损失涨跌幅大小信息，但能减少短期收益噪声对标签的影响，也更符合本文当前的数据条件。

本章使用三类数据：仅基于历史行情的 7 日 K 线监督数据、由 qwen3-max 在 7 日 K 线样本上生成的 CoT 数据，以及包含新闻摘要、研报摘要和 5 日 K 线的评测数据。前两类数据用于监督微调，第三类数据用于观察外部文本加入后，未来 1、3、7 日方向判断是否发生改善。这样安排把训练和评测分开，避免把历史文本覆盖不完整的问题直接带入主训练集。

\section{7 日 K 线监督数据构建}

7 日 K 线监督数据由历史日线行情记录构造。原始数据包含股票代码、交易日期、开盘价、最高价、最低价、收盘价、成交量、成交额、振幅、涨跌幅、涨跌额和换手率等字段。构造样本前，系统先按股票代码分组，再在每只股票内部按交易日期升序排列，得到连续的个股交易序列。

在样本构造时，本文采用滑动窗口方式。对于某只股票的一段连续交易序列，系统取连续 7 个交易日作为输入窗口，记为 Day1 至 Day7；紧接着的第 8 个交易日作为标签日，记为 Day8。若 Day8 的 \texttt{pct\_change} 大于或等于 0，则标签设为“涨”，否则设为“跌”。训练样本默认只保留标签日处于 2021 年 1 月 1 日至 2026 年 1 月 1 日之前的记录，以便和后续 2026 年评测窗口区分。

为了减少不同股票之间价格水平和成交规模差异带来的干扰，系统对每只股票独立估计归一化参数。开盘价、最高价、最低价、收盘价和换手率采用 Min-Max 归一化；成交量先经过 \(\ln(1+x)\) 变换后再进行 Min-Max 归一化；涨跌幅采用 Z-score 标准化；振幅保留原始数值。归一化参数仅基于该股票在训练标签窗口内的数据估计，并在该股票所有训练样本中复用。

最终生成的基础行情数据集为两列结构：\texttt{prompt} 和 \texttt{label}。其中 \texttt{prompt} 包含 7 行归一化 K 线描述和明确任务说明，要求模型只输出“涨”或“跌”；\texttt{label} 保存真实方向标签。当前该数据集共包含 339287 条训练样本，是本文行情方向预测的基础监督数据。

\section{CoT 监督微调数据构建}

直接使用 \texttt{prompt,label} 二列数据可以训练模型输出最终涨跌方向，但不容易约束模型给出可读、稳定的分析过程。因此，本文继续构建 CoT 风格监督微调数据，让模型在输出最终方向之前，先根据 K 线趋势、量能变化和波动特征给出简短推理。这种“先分析、后输出答案”的组织方式与 Chain-of-Thought Prompting 的基本思想一致，即通过显式中间步骤增强模型的分析能力与输出稳定性\cite{wei2022cot}。

CoT 数据通过批量调用 qwen3-max 构建。系统读取基础 7 日 K 线数据集，将每条样本的 \texttt{prompt} 和真实 \texttt{label} 封装为批处理请求。请求中会明确给出当前样本的标准答案，并要求 qwen3-max 在已知正确标签的前提下，用中文生成一段能从行情特征自然推出该答案的简短分析，最后单独一行输出与标准答案一致的“涨”或“跌”。这里并不是让 qwen3-max 重新标注标签，而是利用其生成能力，为已有标签补充解释性训练答案。

批处理流程分为两个阶段。第一阶段，系统生成批处理请求并上传至远程批处理服务，每条请求通过 \texttt{custom\_id} 与原始样本建立一一对应关系。第二阶段，在任务完成后，再根据 \texttt{custom\_id} 将模型生成结果与原始样本对齐。

合并后的 CoT 数据为 \texttt{prompt} 与 \texttt{completion} 两列。新的 \texttt{prompt} 要求模型先在指定标记内分析，再输出最终答案；\texttt{completion} 保存 qwen3-max 生成的推理过程与最终方向标签。当前该数据集共包含 599871 条样本，用于后续 LoRA 监督微调。

\section{多源信息评测数据构建}

为了观察外部文本是否会改善预测效果，本文构建了包含 K 线、新闻摘要和研报摘要的评测数据。该数据的输入包括历史行情记录、新闻摘要、研报摘要以及对应的研报元数据。

在该数据集中，每条样本以目标日期为分析时点。行情部分使用目标日期及之前连续 5 个交易日的归一化 K 线；文本部分检索目标日期之前可见的新闻摘要和研报摘要，并设置最大回看窗口为 30 日。为控制上下文长度，新闻摘要最多保留 30 条，研报摘要最多保留 3 条。新闻和研报正文并不直接进入评测 Prompt，而是先经过摘要压缩后再使用。

标签由未来 1、3、7 个交易日收盘价相对目标日期收盘价的方向组成。若未来收盘价大于或等于目标日期收盘价，则对应位置记为“涨”，否则记为“跌”。因此，每条样本的标签形如“涨，跌，涨”，可以同时评估模型在短期、中短期和一周尺度上的方向判断能力。

该数据集主要用于离线评测，而不是当前 SFT 主训练集。原因是历史时点的文本材料很难完整回放，新闻、研报摘要质量也会受到采集覆盖率和摘要模型影响。本文先用 7 日 K 线 CoT 数据完成较稳定的监督微调，再用这组评测数据观察外部文本接入后的变化。

\section{LoRA 监督微调方法}

\subsection{微调目标}

监督微调阶段主要让基础模型熟悉两件事：一是归一化 K 线特征与涨跌方向标签之间的关系，二是按固定约束生成可解析的回答。对于 CoT 风格数据，模型不仅要输出“涨”或“跌”，还要在指定标记之间给出一段与最终答案一致的分析，以提高回答的可读性和稳定性。

训练目标仍然是自回归语言模型的条件生成目标。给定用户侧 \texttt{prompt}，模型需要生成助手侧 \texttt{completion}。训练时，损失函数只作用于助手回复部分，用户输入部分的标签被设置为 \(-100\) 进行屏蔽。这样可以避免模型把训练目标浪费在复现输入提示词上，而是集中学习如何在给定特征和任务说明后生成符合格式的预测回答。

\subsection{训练数据格式}

监督微调阶段要求训练数据包含 \texttt{prompt} 和 \texttt{completion} 两列。对于每条样本，系统将 \texttt{prompt} 作为用户消息，将 \texttt{completion} 作为助手消息，再通过 Qwen tokenizer 的 chat template 转换为完整对话文本。该做法与指令模型的原始对话格式保持一致，能够减少训练格式与推理格式之间的偏差。

在分词过程中，系统分别构造“仅用户输入”的 prompt token 序列和“用户输入 + 助手回复”的完整 token 序列。完整序列减去 prompt 部分后得到 response token，并将 prompt 部分的 label 全部设为 \(-100\)。若完整序列超过最大长度，则优先保留回复部分，并从 prompt 左侧截断较早的输入 token。当前最大长度默认设为 2048 tokens。

为了提高多卡训练效率，系统还会根据训练数据、模型路径、最大长度和验证集设置构建 tokenized cache。主进程负责将原始 CSV 映射为分词后的 Hugging Face Dataset 并写入缓存，其余进程等待缓存完成后再加载。该机制避免了多进程重复分词，也减少了长时间训练前的数据准备开销。

\subsection{训练参数与实现设置}

本文选择 LoRA 而不是全参数微调，主要考虑到 Qwen2.5-7B-Instruct 的显存和训练成本较高，而本文更关注在既有基础模型上快速适配任务格式\cite{hu2021lora}。训练过程中通过 PEFT 注入 LoRA adapter，覆盖注意力投影层和前馈网络关键线性层。这样做的代价是 adapter 表达能力受秩设置限制，但优点是训练、保存和后续加载都更轻量，也便于在评测时与基础模型对比。

监督微调阶段采用多卡训练，并结合梯度累积、混合精度和梯度检查点控制显存占用。训练过程保留 LoRA adapter 的保存与恢复机制，便于在较长训练过程中继续恢复和比较不同检查点。

训练过程支持从 checkpoint 自动恢复，并保留独立验证环节，用于监控训练损失和选择较优检查点。具体超参数与运行配置已统一整理至附录。

\section{强化学习优化设计}

\subsection{强化学习的引入目的}

监督微调可以让模型学会按指定格式回答，但它本质上仍是在拟合已有样本。对本文任务来说，格式只是第一步，更重要的是预测是否命中、分析是否和结论一致、最终标签能否稳定解析。因此，SFT 之后继续引入强化学习，主要是为了把这些评价要求直接写进训练过程。这与指令对齐研究中在监督微调之后继续利用奖励信号优化输出的思路一致\cite{ouyang2022instructgpt}。

结合当前数据条件，强化学习阶段暂不使用新闻和研报评测样本，而是改用纯 7 日行情窗口构造的多周期方向数据。这样做主要是为了控制变量：文本材料在历史时点上很难完整回放，若直接进入奖励学习，模型出错时很难判断问题来自文本覆盖、摘要质量，还是方向预测本身。行情样本覆盖范围更广，标签也更直接，更适合作为 GRPO 阶段的起点。

本文当前实验主线仍以 SFT 和离线评测结果为主，尚未把 RL 后模型的最终准确率写入结果表。不过，当前系统已经具备可直接运行的多周期方向 GRPO 训练流程，因此强化学习部分不只是概念设想，而是包含数据集、Prompt 与奖励函数的可执行训练方案。

\subsection{GRPO 优化思路}

传统 RLHF 常使用 PPO，但 PPO 通常需要额外训练价值模型，在大模型场景下会增加显存占用和训练复杂度\cite{schulman2017ppo,ouyang2022instructgpt}。本文更倾向采用 GRPO，是因为它不需要额外维护独立 critic。GRPO 在同一输入上采样多个候选输出，再用组内相对得分估计基线，这一思路与 DeepSeekMath 中面向大模型训练效率的设计取向一致\cite{shao2024deepseekmath}。

对于本文任务，同一条 7 日 K 线样本可以生成多组候选预测答案，例如不同的分析过程，以及未来 1 个、3 个、7 个交易日的方向组合。系统随后根据真实标签、输出格式和理由一致性对候选答案打分。得分较高的输出会在组内比较中获得更大优化权重，推动模型逐渐偏向更准确、更稳定、也更符合格式约束的回答。

当前实现中，GRPO 不再沿用 SFT 的单周期 CoT 数据，而是使用纯 7 日行情窗口构造的多周期方向数据作为训练集，并使用 2026 年评测窗口内的同口径样本作为验证集。每条样本仍以过去 7 个交易日的归一化行情描述作为输入，但输出改为先在固定标记内给出分析，再分别输出未来 1 个、3 个、7 个交易日相对 Day7 收盘价的“涨”或“跌”结果。也就是说，GRPO 阶段保留“先分析、后作答”的形式，同时把目标从单一方向扩展到多周期联合判断。

训练器实现沿用 TRL 的 \texttt{GRPOTrainer} 和已有训练基础设施，包括 chat template、左侧截断、自动恢复 checkpoint、随机评估子集以及 vLLM rollout 支持等。这里的强化学习并不是与 SFT 完全割裂的另一条训练线，而是在相近任务格式下继续引入奖励驱动优化。

\subsection{奖励信号与稳定性约束}

当前实现中的核心奖励以未来 1 个、3 个、7 个交易日三个方向标签的平均正确率为主。设模型对某条样本输出的三个预测分别为 \(\hat{y}_1,\hat{y}_3,\hat{y}_7\)，真实标签分别记为 \(y_1,y_3,y_7\)。若模型能够从最终输出中解析出这三个合法的“涨/跌”标签，则奖励定义为三项正确率的平均值，即

\[
r=\frac{1}{3}\Big(\mathbb{I}(\hat{y}_1=y_1)+\mathbb{I}(\hat{y}_3=y_3)+\mathbb{I}(\hat{y}_7=y_7)\Big).
\]

若模型未能按要求输出三项可解析的最终方向标签，则直接记为 0 分。这样定义后，reward 与多周期方向分类目标保持一致，模型在训练中需要同时关注短周期和稍长周期判断，而不是只优化某一个单独预测位置。

日志中还记录五个辅助指标。\texttt{direction\_parse\_success\_rate} 用于统计模型输出能否被成功解析为合法三标签结果；\texttt{direction\_full\_match\_rate} 用于统计三项标签全部命中的比例；\texttt{direction\_1d\_match\_rate}、\texttt{direction\_3d\_match\_rate} 和 \texttt{direction\_7d\_match\_rate} 分别观察三个周期上的表现。这些指标当前不直接写入 reward 标量，但有助于区分两种情况：模型是真的提高了方向判断，还是只是更稳定地按格式作答。

强化学习阶段还需要避免破坏 SFT 已经学到的指令遵循能力。可行做法包括限制每次策略更新幅度，保留相对基础模型或 SFT 模型的 KL 约束，并在奖励中惩罚输出缺失、标签数量错误、过长分析和无法解析等情况。这样可以把优化重点放在预测质量和输出可靠性上，减少模型为追求局部奖励而产生不稳定回答的风险。

在当前训练设置中，稳定性约束主要落在 Prompt、reward 和训练参数上。模型必须先在“思维链开始/结束”标记之间给出分析，再按固定三行格式分别输出未来 1、3、7 个交易日的“涨”或“跌”，为 reward 解析提供稳定入口；若模型未能输出可识别的三项方向标签，该样本自动按 0 分处理；训练过程保留 GRPO 的 \(\beta\) 参数接口，必要时可启用相对参考策略的 KL 正则；上下文长度和 rollout 方式也设置了约束，以维持多卡训练效率和生成稳定性。

\subsection{训练流程与实现设置}

多周期 GRPO 训练的流程包括数据映射、组内采样、奖励计算和策略更新。系统读取 7 日行情多周期数据，将每条 CSV 样本映射为 chat-style Prompt；在每个训练 step 中，对同一输入采样多条候选 completion；系统再从 completion 中解析未来 1、3、7 个交易日的方向标签，与真实值逐项比较，并据此给出平均 reward；最后根据组内相对得分更新当前策略。

在实现层面，该方案延续了基础训练框架中的自动恢复 checkpoint、验证子集评估和 vLLM rollout 支持等机制，以降低多卡训练中的额外开销并提升可重复性。

在部署方式上，GRPO 采用多卡训练与在线 rollout 相结合的方式，以在不额外引入独立推理服务的情况下完成组内采样。若算力受限，也可以改用较轻量的生成方式调试，但训练效率会明显下降。GRPO 训练需要安装完整训练依赖，具体配置见附录。

\section{基础模型、SFT模型与RL模型关系}

本文的模型迭代关系分为基础模型、SFT 模型和 RL 模型三层。基础模型是未经任务微调的 Qwen2.5-7B-Instruct，具备通用语言理解和金融文本分析能力，但尚未针对股票方向预测任务适配输出格式、推理重点和最终标签稳定性。SFT 模型是在基础模型上加载 LoRA adapter 后得到的任务适配模型，通过“输入—输出”样本学习归一化行情特征、预测格式和可解析涨跌方向之间的关系。

RL 模型不是从零开始训练，而是在 SFT 模型已经具备基本任务能力之后，再利用奖励或偏好信号优化。基础模型提供通用语言能力，SFT 模型完成任务格式和领域样本适配，RL 模型则面向最终评价目标继续调整。方法上，RL 阶段既可以直接从基础模型出发，也可以从 SFT 检查点继续训练；后者更符合“先学会按格式做题，再用奖励把题做得更准”的思路。

在评估实现上，本文对基础模型和 SFT LoRA 模型进行对比。评测过程使用 vLLM 加载基础模型，并通过 LoRARequest 动态加载微调 adapter；随后对同一批 eval prompts 分别生成基础模型输出和微调模型输出，再从模型回复中提取最终的“涨”或“跌”标签，与真实标签进行比较。对于多源信息样本，则按照 1、3、7 日三个位置分别计分，计算平均方向预测得分。

\section{本章小结}

本章把训练数据和模型训练流程拆开说明。监督微调部分依赖 7 日 K 线样本和 qwen3-max 生成的 CoT 回答，目标是让模型熟悉行情输入、分析过程和最终方向标签之间的对应关系。多源评测数据则单独保留，用来观察新闻和研报摘要是否带来额外信号。GRPO 部分使用纯 K 线多周期数据，奖励直接来自未来 1、3、7 日方向的逐项命中情况。这样安排后，SFT、多源评测和 RL 优化各自的作用边界更加清楚。
